\begin{helper}
Let \(G\) be a finite loop graph which is an \((n,d,\lambda)\)-graph in the
sense of \(\noderef{LoopGraphNdLambda}\).  Then the loop-complement
\(\noderef{LoopGraphComplement}(G)\) is an \((n,n-d,\lambda)\)-graph.
\end{helper}
\begin{proof}
Write \(\overline G=\noderef{LoopGraphComplement}(G)\).  The vertex set is not
changed by taking the loop-complement, so the vertex count remains \(n\).  If
\(x\overline G y\), then \(xGy\) is false; since \(G\) is symmetric, \(yGx\)
would imply \(xGy\).  Thus \(yGx\) is also false, so \(y\overline G x\).
Hence \(\overline G\) is symmetric.

For a vertex \(v\), the definition \(\noderef{LoopGraphDegree}\) counts the
vertices adjacent to \(v\).  In \(G\) exactly \(d\) vertices are adjacent to
\(v\).  The complement adjacency condition partitions the \(n\) vertices into
the \(d\) vertices adjacent to \(v\) in \(G\) and the vertices adjacent to
\(v\) in \(\overline G\).  Therefore \(v\) has degree \(n-d\) in
\(\overline G\).

It remains to check the non-principal eigenvalue bound from
\(\noderef{LoopGraphNdLambda}\).  Let \(\mu\) be a non-principal adjacency
eigenvalue of \(\overline G\), witnessed as in
\(\noderef{LoopGraphNonprincipalEigenvalue}\) by a nonzero function \(f\) whose
sum over all vertices is zero and for which
\[
  A_{\overline G} f=\mu f.
\]
By \(\noderef{LoopGraphAdjacencyAction}\), for every vertex \(v\),
\[
  (A_{\overline G}f)(v)
  =\sum_{w:\,v\overline G w} f(w)
  =\sum_w f(w)-\sum_{w:\,vGw} f(w)
  =-(A_Gf)(v),
\]
because \(\sum_w f(w)=0\).  Thus \(A_Gf=(-\mu)f\), so \(-\mu\) is a
non-principal adjacency eigenvalue of \(G\).  Since \(G\) is an
\((n,d,\lambda)\)-graph, \(|-\mu|\le\lambda\).  As \(|-\mu|=|\mu|\), the same
bound holds for \(\mu\).  The condition \(\lambda\ge0\) is inherited directly
from the hypothesis that \(G\) is an \((n,d,\lambda)\)-graph.  These four
properties are exactly the assertion that \(\overline G\) is an
\((n,n-d,\lambda)\)-graph.
\end{proof}
